

\fancyhf{}
% \lhead{\begin{tabular}{|p{8.25cm}} \hline {\large \bf SNO+ Laser Procedures} \\ \\ \\ \\ \end{tabular}}
\lhead{\begin{tabular}{|p{8.25cm}|p{6cm}|}  
	\hline 
		{\large \bf Salt Oven Usage Procedures}  
										& Document No:  \\ 
								\cline{2-2} & Revision No: 01\\  
								\cline{2-2} & Effective Date: 2018-04-30\\ 
								\cline{2-2} & Page \thepage \\  \end{tabular}}




\begin{tabular}{||l|l|l||}
\hline\hline
& \multicolumn{2}{p{8cm}||}{\bf Salt Oven Usage Procedures} \\
\includegraphics[width=6cm]{figures/SNOplus_logo.png} & \multicolumn{2}{p{8cm}||}{} \\
\hline
\multicolumn{2}{||p{8.5cm}|}{Document Number:} & Revision Number: 01\\
\hline
\multicolumn{3}{||l||}{Document Owner: SNO+ Calibration Post-Doc} \\
\hline
\multicolumn{3}{||l||}{Reviewer:}\\
\hline
Name: & Signature & Date \\
\hline
\multicolumn{3}{||l||}{Authorizer:}\\
\hline
Name: & Signature & Date \\
\hline\hline
\end{tabular}
\thispagestyle{empty}


\section{Purpose}
This document describes the procedure for preparing and operating the Barnstead Thermolyne F21100 for use with a Sodium Chloride (NaCl) loaded quartz tube. The tube furnace heats the salt to a desired temperature of ~610C, just below its melting point. A stream of Helium flows through the heated tube, along with the heated NaCl aerosol. The tube has a water cooling section before the gas outlet, preventing damage to the following gas lines and the Umbilical. 

These procedures are meant to suppliment the use of the DT generator
in the operation of the Cherenkov source for the production of
$^{8}$Li. Any operations that will require the salt oven will also
require the DT Generator procedures. Some operations given here will
supercede those given in the DT Generator procedures as stated in the
Cherenkov deployment procedures. 

\paragraph{Warnings}
\begin{small}
Refer to THA for required precautions needed for tube furnace operations. 
Always use heat resistance gloves when handling the salt tube.
Ensure a proper barricade is in place around the oven work space.
Ensure proper signage is in place to warn others of the hazards present beyond the barricade. 
\end{small}

\section{Procedures}
\subsection{Preparing Oven}
\begin{enumerate}
\item \CheckBox[name=oprep1]{} Only use Tube Furnace in designated area
\item \CheckBox[name=oprep2]{} Make sure the furnace is empty before loading the salt tube
\item \CheckBox[name=oprep3]{} Load the salt tube in the oven
\item \CheckBox[name=oprep4]{} Connect Helium gas inlet and outlet to the tube
\item \CheckBox[name=oprep5]{} Connect the output end to the water cooling loop using the attached flex hose
\end{enumerate}

\subsection{Starting Water Cooling Loop}
\begin{enumerate}
\item \CheckBox[name=cloop1]{} Ensure the needle valve behind the gas board is closed.
\item \CheckBox[name=cloop2]{} Check that the inlet line is attached to the water line after the valve and that the outlet line is running into the drain channel.
\item \CheckBox[name=cloop3]{} Insert the input line into the sleeve on the salt tube closer to the oven. Clamp the line using the hose clamp provided.
\item \CheckBox[name=cloop4]{} Insert the output line into the sleeve on the salt oven nearer to the He gas outlet. Clamp the line using the hose clamp provided.
\item \CheckBox[name=cloop5]{} Slowly open the valve. Ensure that all of the bubbles are out of the line before starting the Helium flow.
\item \CheckBox[name=cloop6]{} Check the seal on the input and output lines for water leaks. If the water pressure is too high water will leak from the cooling loop connections. If there is a leak close the water valve until the leak is no longer present.
\end{enumerate}


\subsection{Turn on Oven}
\begin{enumerate}
\item \CheckBox[name=toov1]{} {Turn on the furnace, make sure the setpoint is 615C}
\item \CheckBox[name=toov2]{} {Use the “UP” or “DOWN” button while in the home display to change the setpoint}
%\end{eu}
\subsubsection{Tuning Cycle}
%\begin{itemize}
\item \CheckBox[name=toov3]{} {Press the “PAGE” button until “Atun” is displayed}
\item \CheckBox[name=toov4]{} {Press the “SCROLL” button, the display will read \textbf{tunE}}
\item \CheckBox[name=toov5]{} {Press either “UP” or “DOWN”  to select \textbf{ON}}
\item \CheckBox[name=toov6]{} {Press both “SCROLL” and “PAGE” to return to the main display }
\item \CheckBox[name=toov7]{} {Wait for the tuning to complete}
%\end{}
\subsubsection{Continuous Heating}
%\begin{itemize}
\item \CheckBox[name=toov8]{} {Press the “PAGE” button until \textbf{SP} is displayed }
\item \CheckBox[name=toov9]{} {Press the “SCROLL” button until \textbf{tm.OP} is displayed }
\item \CheckBox[name=toov10]{} {Press either “UP” or “DOWN”  to select \textbf{Opt.2}. This setting will ramp and maintain setpoint temperature}
\item \CheckBox[name=toov11]{} {Press the “SCROLL” button until \textbf{StAt} is displayed}
\item \CheckBox[name=toov12]{} {Press either “UP” or “DOWN”  to select \textbf{run} }
\item \CheckBox[name=toov13]{} {Press both “SCROLL” and “PAGE” to return to the main display  }
\end{enumerate}

\section{Flow Gas}
\begin{enumerate}
\item \CheckBox[name=gflow1]{} Ensure that all of the valves on the gas board are closed.
\item \CheckBox[name=gflow2]{} Check that the valves on the He, N$_{2}$, and CO$_{2}$ bottles are closed.
\item \CheckBox[name=gflow3]{} Verify that the "Solenoid Valve II" switch is in the OFF position.
\item \CheckBox[name=gflow4]{} Enter the DCR using standard entry procedure and verify that
	\begin{itemize}
	\item \CheckBox[name=gflow5]{} the gas input line is connected to the "dry" end of the umbilical
	\item \CheckBox[name=gflow6]{} the gas return line is connected to the "dry" end of the umbilical.
	% \item the pressure transducer is powered up and that a voltmeter is connected to it
	\end{itemize}
\item \CheckBox[name=gflow7]{} Turn on the power onf the main instrument panel and the dual flow controller box.
\item \CheckBox[name=gflow8]{} Flip the switch on the flow controller box to "B" (for N$_{2}$ and/or CO$_{2}$). The reading takes a few minutes to equilibrate to zero or near zero. 
\item \CheckBox[name=gflow9]{} Ensure that the CO$_{2}$ and He bottles are hooked up to the input line at VA2. Ensure that the N$_{2}$ bottle is connected to VA1.
\item \CheckBox[name=gflow10]{} Check that VA2 and VA2 are closed.
\item \CheckBox[name=gflow11]{} Close the needle valve on the CO$_2$ bottle.
\item \CheckBox[name=gflow12]{} Open the main valves on the CO$_{2}$, He and N$_{2}$ bottles.
\item \CheckBox[name=gflow13]{} Set the regulator on the CO$_{2}$ and He bottle to 33~psi. Set the N$_{2}$ regulator to 20~psi.
\item \CheckBox[name=gflow14]{} Record the bottle pressure:
\begin{center}
\begin{tabular}{|c|c|}
\hline
Transducer & Reading \\
\hline
CO$_{2}$ Bottle Pressure & \\
\hline
He Bottle Pressure & \\
\hline
N$_{2}$ Bottle Pressure & \\
\hline
\end{tabular}
\end{center}
{\emph The pressure of a full CO$_{2}$ bottle is 850 to 900~psi. The
  CO$_{2}$ in the bottle is in liquid or solid form and the pressure
  will stay relatively high until gas only and will drop rapidly
  thereafter. The helium bottle is simply gas under pressure and the
  pressure will drain steadily throughout its use. }
% \item Open the inlie needle valve after the CO$_{2}$ regulator.
\item \CheckBox[name=gflow15]{} Open SVII by setting the switch on the main control panel to
  {\bf MAN}.
\item \CheckBox[name=gflow16]{} Open the manual shut off valve VB2
\item \CheckBox[name=gflow17]{} Direct flow to the He flow meter using valve VD1.
\item \CheckBox[name=gflow18]{} Set the He flow meter to manual and turn the control knob one
  turn so that some gas will flow.
\item \CheckBox[name=gflow19]{} Direct flow towards the He flow meter using valve VD1.
\item Direct flow to the Salt Oven using valve VF1.
\item \CheckBox[name=gflow20]{} Open SVI by setting the switch on the main control panel to
  {\bf MAN}.
\item \CheckBox[name=gflow21]{} Open the manual shut off valve VB1.
\item \CheckBox[name=gflow22]{} Set the N$_{2}$ flow meter to manual and turn the control knob
  one quarter turn.
\item \CheckBox[name=gflow23]{} Direct the N$_{2}$ flow toward P1 (and VF1) so that the N$_{2}$
  mixes with the He flow.
\item \CheckBox[name=gflow24]{} Open valve VF2 to accept the returning CO$_{2}$.
\item \CheckBox[name=gflow25]{} Open valve VC2.
\item \CheckBox[name=gflow26]{} Direct the flow towards the exhaust using valve VB4.
\item \CheckBox[name=gflow27]{} Open the exhast needle valve VA4 fully.
\item \CheckBox[name=gflow28]{} {\bf Slowly} open valve VA2 and look to see if gas is flowing.
\item \CheckBox[name=gflow29]{} {\bf Slowly} increase the flow until you reach around a 200-280
  reading. Make sure P1 never goes above 100 PSI. Remember that the
  time constants are long.
\item \CheckBox[name=gflow30]{} Record gas Pressures and Flow.
\begin{center}
\begin{tabular}{|c|c|}
\hline
Transducer & Reading \\
\hline
P1 & \\
(PSIA) & \\
\hline
P2 & \\
(PSIA) & \\
\hline
P1 & \\
(PSIA) & \\
\hline
P2 & \\
(PSIA) & \\
\hline
Flow & \\
(cc/s $\times$ (1/1000)) & \\
\hline
Target & \\
Position & \\
\end{tabular}
\end{center}
  
\end{enumerate}

\section{Turn off Oven}
\begin{enumerate}
\item \CheckBox[name=ofov1]{} {Press the “PAGE” button until \textbf{SP} is displayed }
\item \CheckBox[name=ofov2]{} {Press the “SCROLL” button until \textbf{StAt} is displayed}
\item \CheckBox[name=ofov3]{} {Press either “UP” or “DOWN”  to select \textbf{off} }
\item \CheckBox[name=ofov4]{} {Press both “SCROLL” and “PAGE” to return to the main display  }
\item \CheckBox[name=ofov5]{} {The oven should start to lose heat, let the salt tube cool to ambient temperature}
\item \CheckBox[name=ofov6]{} {Turn off tube furnace}
\end{enumerate}


